% ID: SEL-D0064-OA-E03
% ESTADO: validada. S064.
Sean $C$ y $M$ los sucesos de leer \textit{CIUDAD} y \textit{LA MAÑANA}, respectivamente.

\textbf{a)} Por complemento, $1-P(C\cup M)=1-0{,}85=\boxed{0{,}15}$.

\textbf{b)} La unión se divide en quienes leen $C$ y quienes leen sólo $M$. Por tanto, $P(M\cap C^c)=0{,}85-0{,}70=\boxed{0{,}15}$.

\textbf{c)} $P(M)=0{,}15+0{,}18=0{,}33$ y $P(C\cap M^c)=0{,}70-0{,}18=0{,}52$. Entre quienes no leen $M$, la proporción que lee $C$ es
\[
P(C\mid M^c)=\frac{0{,}52}{0{,}67}=\boxed{\frac{52}{67}}\approx0{,}7761.
\]
